\reporttitle{第三阶段协同仿真与性能分析文档}


工程组合为0f＋firstdot V2＋原C8。本阶段围绕架构一致性、执行周期、软件调用和实现成本展开评估。\par


\section{验证与评估目标}


本阶段以独立NEMU模型与香山RTL进行差分验证，检查完整向量结果、异常和控制行为；随后通过同条件RVV对照测量执行周期，结合C8调用与综合时序结果评估工程效果。\par


\section{NEMU扩展}





执行过程先匹配custom-0编码并检查合法向量状态，再快照读取两源，逐项转换signed int8、从零求和并写回完整128位结果。独立软件求和有助于与硬件乘法/归约结构交叉验证。\par


平台适配对齐非法指令xTVAL行为。固定香山中，译码已知非法与FU发现VL非法的路径具有不同xTVAL表现，reference按同一架构谓词建模。差分检查覆盖异常码、陷阱PC、xTVAL及目的状态，全核实验在M模式下执行。\par


\section{差分、拒错与恢复}


DiffTest将硬件提交后的架构状态与NEMU对应状态比较。除正常算术结果外，本项目检查完整向量写回、向量CSR、寄存器别名和非法状态路径。\par


负控对vdot结果的bit100注入错误，检验高96位比较是否真正生效。图S3-03展示2026年10月24日原0f记录中的negative-vdot拒错与functional-0同镜像恢复两次运行，绑定EMU 2d1db77c…与BIN b8b7e1a7…。原高64位的第36位对应完整向量的bit100；首次失配PC为0x80001bdc，机器字0x0294050b，目的v10。恢复正常reference后，15个计划用例全部通过并以GOOD TRAP、exit 0结束。\par


后续R5/G5在另一组冻结运行身份下选择原13项记录与1项新控制诊断组成验收，原始记录重新解析后通过。全核功能矩阵包含117个不同用例及124次分片执行。新控制诊断检查131072字节寄存器结果。各层测试分别记录输入向量、事务、用例和执行次数。\par




R5 clean-02原宿主UTC并发审计存在80.119ms表观四路重叠，旧FAIL记录保留。补证以同boot/PID单调包络完成所选审计，组合UTC峰3、补项单调峰1。系统UTC稳定性仍待建立，性能采用DUT guest周期。\par


\section{等价RVV基线与测量方法}


已知配置的RVV单次点积通过vwmul.vv得到宽化乘积，再切换到e16、清零目的向量，使用vwredsum.vs形成32位总和，最后恢复入口向量状态。目的清零覆盖完整向量，宽化归约保留已清零的高位。RVV批处理共享允许共享的配置开销。\par


两路径计算相同输入和完整输出，检查声明的寄存器副作用、CSR以及缓冲区护栏。相同镜像使用相同窗口规则、数据和循环次数。正式周期按三组数据、每组四个正式窗口统计，首次及预热另存，独立重放用于一致性检查。\par


加速比 $S = \frac{C_{\mathrm{RVV}}}{C_{\mathrm{vdot}}}$。相对旧版的周期降低率 $R = \frac{C_{\mathrm{old}} - C_{\mathrm{new}}}{C_{\mathrm{old}}}$。\par


在固定香山实现中，fence与CSR读组成的测量边界经实际入队、采样、退休记录核验。每窗校验四条CSR指令的PC、机器字和ROB身份，cycle差与全核tick差对应，instret差与实际退休/ELF展开计数一致。绝对计数器按整数解析。原记录中的首次、预热和失败样本完整保存。\par


\section{硬件微基准结果}


下表来自0f同镜像对照合同，均为正式窗口均值。\par


{\small\setlength{\tabcolsep}{3pt}\begin{longtable}{>{\raggedright\arraybackslash}p{0.25\linewidth}>{\raggedright\arraybackslash}p{0.21\linewidth}>{\raggedright\arraybackslash}p{0.21\linewidth}>{\raggedright\arraybackslash}p{0.23\linewidth}}
\toprule
\textbf{场景} & \textbf{vdot cycles} & \textbf{RVV cycles} & \textbf{加速比}\\ \midrule\endhead
单次点积 & 71.000 & 85.250 & 1.200704×\\
依赖链256 & 1857.000 & 5440.500 & 2.929725×\\
四路独立batch128 & 221.000 & 678.000 & 3.067873×\\
加载/计算/存储stream64 & 1127.500 & 1133.000 & 1.004878×\\
\bottomrule\end{longtable}}


独立批量可以利用O2的连续接收能力，依赖链受到前后数据依赖约束，stream场景还承担访存及循环开销。结果显示收益随工作负载改变。\par


O1→O2的独立同镜像比较中，batch128由362.5降至221周期，减少39.0345\%。H1及0f修复分别保持各自96对正式窗口周期相同。阶段间比较采用各自冻结合同。\par




\section{软件调用与C8结果}


firstdot以首段点积初始化累加状态，后续段按约定求和。C8在已验证8×32路径中复用配置，提供普通C显式入口和RVV回退。输入为对齐的288字节A||x，输出为8个16字节完整槽。\par


{\small\setlength{\tabcolsep}{3pt}\begin{longtable}{>{\raggedright\arraybackslash}p{0.25\linewidth}>{\raggedright\arraybackslash}p{0.69\linewidth}}
\toprule
\textbf{同合同8×32实现} & \textbf{完整调用cycles}\\ \midrule\endhead
firstdot & 282.25\\
原C8配置复用vdot & 273.00\\
本批最快合格RVV & 299.00\\
\bottomrule\end{longtable}}


C8相对firstdot减少3.2772\%，相对同批最快RVV加速1.095238倍。三数据集完整矩阵给出周期比较，普通C调用、软件包与最终镜像检查验证了ABI、回退和接入方式。\par


原合格镜像、复制对照、SDK接入版BIN均为31440字节，SHA相同。PT\_LOAD、48条调用表项与12组输入数据一致，ELF差异位于调试节。627项输入与封存材料的身份核对，将SDK接入结果与原功能记录对应起来。\par


4×8和16×64采用原firstdot入口，4×13选择更快的RVV。各尺寸沿用对应的已验证调用入口。\par


\section{面积、时序与设计成本}


使用Nangate45 Typical 1.10V/25°C，Yosys0.9、ABC1.01及OpenSTA2.0.17。主档周期1ns、负载5fF，采用固定线负载模型，另有0.5ns、2ns及1fF、20fF敏感性档。\par


{\small\setlength{\tabcolsep}{3pt}\begin{longtable}{>{\raggedright\arraybackslash}p{0.27\linewidth}>{\raggedright\arraybackslash}p{0.3\linewidth}>{\raggedright\arraybackslash}p{0.35\linewidth}}
\toprule
\textbf{指标} & \textbf{FU} & \textbf{完整VFEX4父级}\\ \midrule\endhead
标准单元面积 μm² & 17,298.246 & 284,614.946\\
nominal r2r到达 ns & 1.163404 & 3.261218\\
主档i2r setup裕量 ns & −0.268520 & −2.393998\\
\bottomrule\end{longtable}}


父级面积包含FU及DIV。负裕量表明1ns参考目标仍有时序违例，整核Fmax、功耗和布线后时序为待测。H1相对O2独立FU面积降低2.33386\%，父级降低0.137349\%，同时保留部分路径退化。0f的取消关钟修复以正确性为依据，其周期不退化与成本变化分别留档。\par


\section{优化探索及最终取舍}


{\small\setlength{\tabcolsep}{3pt}\begin{longtable}{>{\raggedright\arraybackslash}p{0.27\linewidth}>{\raggedright\arraybackslash}p{0.3\linewidth}>{\raggedright\arraybackslash}p{0.35\linewidth}}
\toprule
\textbf{探索} & \textbf{结果} & \textbf{最终处理}\\ \midrule\endhead
输入分离候选（P1/P1-R1） & 对照窗口偏差超过预设1\%限值 & 暂缓\\
单拍两槽候选（P4） & 47项成本检查中6项失败 & 未采用\\
单拍两槽修订（P4-R1） & 47项成本检查中10项失败 & 未采用\\
C8分段归约 & 实测273.75→273周期，一致性检查与3\%收益目标未通过 & 未采用\\
FP32到QKV探索（P3） & 未取得完整计时窗口，尚未确认完整调用能否在7200秒预算内完成 & 暂缓\\
\bottomrule\end{longtable}}


综合这些结果，工程选用0f＋firstdot V2＋原C8。应用评估包括点积微基准和固定GEMV调用；另有INT16独立单元扩展实验，生产指令保持INT8。\par


\section{结论}


本项目完成自定义指令硬件接入、独立NEMU扩展、分层功能验证及等价RVV周期分析。O2提升独立吞吐，H1降低结果状态成本，0f处理DIV取消关钟，firstdot与C8改善固定软件调用。固定源、镜像和样本支持结果复算，1ns时序收敛及完整物理评估构成后续开发方向。\par


\href{https://github.com/cc77-builder/riscv/blob/eca41d7cd959f5f3c4f240e382919c492d1ccb96/team/patches/nemu/README.md}{NEMU补丁说明} · \href{https://github.com/cc77-builder/riscv/blob/eca41d7cd959f5f3c4f240e382919c492d1ccb96/team/spec/measurement.md}{测量合同} · \href{https://github.com/cc77-builder/riscv/blob/eca41d7cd959f5f3c4f240e382919c492d1ccb96/team/optimizations/o2/r3-div-gate/PERFORMANCE\_COST\_RESULT.md}{0f性能与成本} · \href{https://github.com/cc77-builder/riscv/blob/eca41d7cd959f5f3c4f240e382919c492d1ccb96/team/docs/version-comparison-r5-g5-20261027.md}{版本数据对照}\par
\section{工程版本与主要产物}
测量软件以T13实现为基础，原文件及SHA-256保存在\seqsplit{team/references/zufan-t13/}，对应提交为\seqsplit{f090665130107282be0990a7664aebf26e491b87}。测量实验按自己的构建、镜像和运行记录建立样本集合。\par
该测量硬件源码提交为\seqsplit{4460300e9557ad20ac152fa67b1827c5d62802dd}，工作负载的构建提交另行记录。\par

NEMU以\seqsplit{a48a11f208ff81ad765591d83271d1b339fb3861}为基点，使用vdot补丁1.1.0及\seqsplit{xs-16ae9dd-vdot-runtime-tval-v1}平台配置。补丁绑定7个变化文件，SHA256为\seqsplit{6c165bf5f94ef0bfb4f894a5f06eb6d620d4a7e0e12f3b1400fa72811f5d24db}。\par

\begin{itemize}
\item 本章引用的历史资料发布提交为eca41d7cd959f5f3c4f240e382919c492d1ccb96。
\item 生产硬件提交为0f12df675378556670afd0dd2153f47fcb330610。
\item clean-02 EMU的SHA256为\seqsplit{acf672eb39f836da4000b253849e3214089e9c8e5198060a57d2cc6bbd393da9}。
\item reference的SHA256为\seqsplit{62b1707c52a4c323468481e0c1863b79f244d9a3fd54aeb7b7d2ecb5f01a35ca}。
\item C8最终BIN的SHA256为\seqsplit{56805506eb98b68269c73a3645d8f75dff73b697e30dec5f93ccc4716089ad34}。
\end{itemize}
以上标识对应本章采用的工程组合及主要运行产物，分项实验使用各自记录的构建与镜像身份。
